\chapter{Encapsulation, Ownership and Concurrency}
\label{ch:encap}

Four Ada features let an interface control, in the type system itself, what client
code is allowed to do, and together they make a peripheral driver that is
misuse-resistant by construction. \texttt{private} hides a type's representation;
\texttt{limited} forbids copying it; \texttt{controlled} gives it automatic setup
and teardown; \texttt{protected} serialises access to it. None costs more at run
time than the discipline you would otherwise have to write and remember by hand,
and each turns a class of ``be careful to\ldots'' rules into a compiler-enforced
guarantee. This chapter explains each keyword with small examples and its embedded
payoff; the \emph{Writing a Task-Safe Peripheral Driver} chapter shows them
combined in the real HAL.

\section{\texttt{private} types: hiding the representation}\index{private types}

A private type is declared in two halves. The visible part of a package gives the
type a name and the operations on it; the \emph{private part} (everything after
the word \texttt{private}) gives its real definition. Client code may use the
name and call the operations, but cannot see or touch the components.

\begin{lstlisting}
package Counters is
   type Counter is private;             -- clients see only this and the ops
   procedure Bump  (C : in out Counter);
   function  Value (C : Counter) return Natural;
private
   type Counter is record               -- the real definition, hidden
      N : Natural := 0;
   end record;
end Counters;
\end{lstlisting}

A user can declare a \texttt{Counter}, \texttt{Bump} it and read its
\texttt{Value}, but cannot write \texttt{C.N}, cannot rely on it being a record,
and will not break if the representation later changes. This is the difference
between a \texttt{struct} in a public C header (whose fields anyone may scribble
on) and a genuinely opaque handle.

\textbf{Embedded payoff.} The representation of a driver handle (which host it
owns, a DMA-descriptor pointer, a hardware index) is hidden, so application code
cannot corrupt it and the driver author stays free to change it.

\section{\texttt{private} child packages: sealing the unsafe layer}

Visibility extends to whole packages. A child package declared \texttt{private} is
visible only inside its own family's bodies; application code cannot even
\emph{name} it in a \texttt{with} clause. The HAL uses this to seal the only code
that touches registers:

\begin{lstlisting}
private package ESP32S3.SPI.Engine is   -- un-with-able from an application
   --  the raw, lock-free register driver: the ONLY code that pokes SPI.
   ...
end ESP32S3.SPI.Engine;
\end{lstlisting}

The public \texttt{ESP32S3.SPI} package \texttt{with}s the engine in its body and
mediates every operation. \textbf{Embedded payoff.} ``Poke the registers
directly'' is not merely discouraged: it is \emph{unwritable}. A
\texttt{with ESP32S3.SPI.Engine;} in an application fails to compile, so the unsafe
core is reachable only through the safe fa\c{c}ade.

\section{\texttt{limited} types: values that cannot be copied}\index{limited type}

A limited type has no assignment (\texttt{:=}) and no predefined equality
(\texttt{=}). You can declare one, pass it (by reference) and operate on it, but
you cannot copy it. Make a type limited when a value stands for a \emph{unique}
thing that must not be duplicated: a handle to a single hardware resource, a
lock, an open session:

\begin{lstlisting}
   type Session is limited private;     -- a SPI bus session
   ...
   A, B : Session;
   B := A;       -- COMPILE ERROR: Session is limited, it cannot be copied
\end{lstlisting}

If you \emph{could} copy a handle, two copies would each believe they own the bus;
releasing one would leave the other dangling. \texttt{limited} makes that bug
impossible to write. It is the real shape of the SPI \texttt{Session}, the GDMA
\texttt{Channel}, the ADC \texttt{Reader}, the LEDC/PCNT \texttt{Channel} and the
rest of the HAL's ownership handles.

\section{\texttt{controlled} types and RAII: deterministic teardown}\index{controlled type}\index{RAII}

\texttt{Ada.Finalization.Controlled} (and its non-copyable sibling
\texttt{Limited\_Controlled}) give a type three hooks the runtime calls
\emph{automatically}: \texttt{Initialize} when an object comes into being,
\texttt{Finalize} when it leaves scope (\emph{including} when an exception unwinds
the scope), and \texttt{Adjust} after a copy (for the non-limited kind). That is
RAII: a resource's lifetime is tied to an object's scope.

\begin{lstlisting}
   type Session is new Ada.Finalization.Limited_Controlled with record
      Host : ...;
   end record;
   overriding procedure Finalize (S : in out Session);   -- auto-release the bus
\end{lstlisting}

\begin{lstlisting}
   declare
      S : Session;                  -- Initialize
   begin
      Acquire (S, Host => SPI2);    -- take the bus
      Transfer (S, ...);
   end;                             -- Finalize runs HERE: the bus is released,
                                    -- even if Transfer raised an exception
\end{lstlisting}

\textbf{Embedded payoff.} You \emph{cannot forget} to release a DMA channel,
unlock a bus, or power down a peripheral: the release is bound to the scope and
happens on every exit path, with no leak, deterministic timing, and no
\texttt{goto cleanup} ladder. \texttt{limited} + \texttt{controlled} is the
project's standard handle: limited stops you copying the owner, controlled
guarantees it is released. (The real SPI \texttt{Session}'s overriding
\texttt{Finalize} is commented ``auto-release on scope exit'' for exactly this.)

\cnote{ESP-IDF hands you an opaque handle and a matching \texttt{*\_delete} /
\texttt{*\_free} you must remember to call on every path (including the error
ones). A controlled \texttt{Session} releases itself when it leaves scope (even
when an exception unwinds past it) so the leak and the missed-cleanup bug
cannot happen.}

\section{\texttt{protected} types and objects: race-free sharing}\index{protected object}

A protected object is data plus operations with mutual exclusion built in: a
monitor. Only one task at a time runs a protected \emph{procedure} or
\emph{entry}; the compiler generates the locking, and on this runtime the lock is
a \emph{priority ceiling}, so worst-case blocking is bounded (real-time
friendly). There are three kinds of operation:

\begin{itemize}
  \item a protected \textbf{function} reads the state and may run concurrently
        with other functions (a shared read);
  \item a protected \textbf{procedure} reads and writes with exclusive access
        (one caller at a time);
  \item a protected \textbf{entry} is a procedure with a Boolean \emph{barrier};
        a caller blocks until the barrier is open (condition synchronisation).
\end{itemize}

\begin{lstlisting}
   protected Counter is
      procedure Bump;                  -- exclusive write
      function  Value return Natural;  -- shared read
   private
      N : Natural := 0;
   end Counter;

   protected body Counter is
      procedure Bump is begin N := N + 1; end Bump;
      function  Value return Natural is (N);
   end Counter;
\end{lstlisting}

\texttt{Counter.Bump} from one task and \texttt{Counter.Value} from another never
race: no \texttt{Volatile}, no hand-rolled critical section. The HAL uses
exactly this to arbitrate a shared peripheral: a per-port \texttt{protected Guard}
(or \texttt{Pool}, or \texttt{Lock}) hands out ownership so two tasks cannot drive
the same controller at once, and an interrupt handler is itself a protected
procedure (the Interrupts chapter). \textbf{Embedded payoff.} Shared state across
tasks, cores and interrupt handlers is correct by construction, with bounded,
analysable blocking.

\section{The four together: a misuse-proof driver}

Combined, these keywords make a peripheral that is hard to misuse even on purpose:

\begin{itemize}
  \item the registers live in a \texttt{private} child: you cannot see or poke
        them;
  \item the public handle is \texttt{limited}: you cannot copy ownership;
  \item it is \texttt{controlled}: you cannot forget to release it;
  \item access is \texttt{protected}: you cannot race it.
\end{itemize}

That is precisely the HAL's task-safe driver skeleton; the \emph{Writing a
Task-Safe Peripheral Driver} chapter builds one end to end (Engine $+$ Session
$+$ Guard). The payoff is that a whole class of embedded bugs (a dangling
handle, a forgotten unlock, a double-driven bus, a torn shared counter) becomes
a compile error or a structural impossibility instead of a field failure.

\section{Related keywords}

A few neighbours round out the family. Each is shown with a small example.

\paragraph{\texttt{constant}: a value that never changes.} A constant is given
its value once and can never be assigned again; the compiler may place it in
read-only memory (flash), saving RAM. A constant with no type (a \emph{named
number}) is a pure compile-time value.

\begin{lstlisting}
   Max_Pins : constant := 48;                 -- named number (no type, no storage)
   Sine     : constant array (0 .. 255) of Integer := ( ... );  -- table, in flash
   --  Sine (K) := 0;  -- COMPILE ERROR: Sine is constant
\end{lstlisting}

In embedded code lookup tables, calibration data and configuration are
\texttt{constant} so they live in flash, not in scarce RAM.

\paragraph{\texttt{aliased}: an object you may point at.} By default you cannot
take \texttt{'Access} of a declared object (the accessibility rules guard against
dangling pointers). Marking it \texttt{aliased} permits it, needed when, say, a
DMA buffer must be handed to a descriptor by reference.

\begin{lstlisting}
   Buf  : aliased DMA_Buffer;          -- 'Access may now be taken
   Desc.Src := Buf'Access;             -- give the DMA engine a pointer to it
\end{lstlisting}

\paragraph{\texttt{renames}: a new name for an existing thing.} A renaming
declaration introduces an \emph{alias}: a second name that denotes the \emph{same}
entity, with no copy and no new object made. It works for objects, components,
subprograms, packages and exceptions.

\begin{lstlisting}
   --  alias a deeply-nested register field: short name, same storage
   Ctrl : Unsigned_32 renames Periph.SPI2.User.Control;
   Ctrl := 16#1#;                      -- writes Periph.SPI2.User.Control

   --  alias a package to shorten a long path (just this one name, unlike `use`)
   package GE renames ESP32S3.GPIO;
   GE.Toggle (2);

   --  alias a subprogram -- give an operation a local or operator name
   function "+" (L, R : Vector) return Vector renames Vectors.Add;
\end{lstlisting}

For an object or component the alias shares the original's storage (writing
through \texttt{Ctrl} above writes the real register field), so \texttt{renames}
is the way to name one specific deeply-qualified thing once and use the short name
afterwards. Note how it differs from \texttt{use} (\S\ref{ch:registers}):
\texttt{use} makes \emph{all} of a package's names directly visible, whereas
\texttt{renames} introduces exactly \emph{one} new name and says precisely what it
stands for, so it shortens a single long reference without the wholesale
visibility (and possible name clashes) that \texttt{use} brings.

\paragraph{\texttt{tagged}: records with inheritance and dispatch.} A tagged
type can be \emph{extended} (a derived type adds components) and its primitive
operations can \emph{dispatch} at run time: a call through a class-wide value
\texttt{T'Class} selects the right override for the actual object.

\begin{lstlisting}
   type Sensor is tagged record Id : Natural; end record;
   function Read (S : Sensor) return Integer;

   type Temp_Sensor is new Sensor with record Pin : Pin_Id; end record;  -- extends
   overriding function Read (S : Temp_Sensor) return Integer;            -- override
\end{lstlisting}

A call \texttt{Read (S)} where \texttt{S} is \texttt{Sensor'Class} dispatches to
\texttt{Temp\_Sensor}'s \texttt{Read} if that is the actual type. It is used
sparingly in embedded code (the \texttt{embedded}-profile demo shows library-level
tagged dispatch) because a dispatching call is an indirect call.

\paragraph{\texttt{abstract}: a type or operation with no implementation.} An
abstract type cannot be instantiated, and an abstract operation has no body:
every concrete descendant \emph{must} provide one. This is how Ada expresses an
interface: a contract that says ``every sensor shall have a \texttt{Read}.''

\begin{lstlisting}
   type Sensor is abstract tagged null record;
   function Read (S : Sensor) return Integer is abstract;  -- no body; must override
   --  S : Sensor;  -- COMPILE ERROR: Sensor is abstract
\end{lstlisting}

These are reached for less often than \texttt{private} /\texttt{limited} /
\texttt{protected} /\texttt{controlled}, but knowing them completes the picture of
how Ada lets the \emph{declaration} of a type state what may -- and may not -- be
done with it.
